% 第6章 6.7.2：七点图上RatioCut阈值扫描的逐步执行（由 generate.py 生成并用分数复算）。
% 两个单位权三角形 {1,2,3}、{4,5,6} 由单位权桥 3-4 相连，顶点 7 以权 0.6 挂在顶点 2 上。
% 按第二特征向量的坐标依次移入 7,2,1,3,4（5 与 6 坐标相同，留在最后一组）。
%   (1) 移入 7：S=[7]，割权 0.6，RatioCut=7/10≈0.7000
%   (2) 移入 2：S=[2, 7]，割权 2，RatioCut=7/5≈1.4000
%   (3) 移入 1：S=[1, 2, 7]，割权 2，RatioCut=7/6≈1.1667
%   (4) 移入 3：S=[1, 2, 3, 7]，割权 1，RatioCut=7/12≈0.5833
%   (5) 移入 4：S=[1, 2, 3, 4, 7]，割权 2，RatioCut=7/5≈1.4000
% 蓝色：S；绿色：S 的补；红色粗边：跨组边；蓝色粗描边：本步移入的顶点；红字：到该步为止的最优值。
\begin{tikzpicture}[x=1mm,y=1mm,font=\figurefont\footnotesize,text=FigureInk,
  vS/.style={circle,draw=FigureBlue,fill=FigureBlueFill,line width=.9pt,minimum size=5mm,inner sep=0pt,font=\figurefont\fontsize{7.5}{9}\selectfont},
  vC/.style={circle,draw=FigureGreen,fill=FigureGreenFill,line width=.9pt,minimum size=5mm,inner sep=0pt,font=\figurefont\fontsize{7.5}{9}\selectfont},
  vN/.style={vS,draw=FigureBlue,line width=1.9pt},
  edge/.style={draw=FigureStroke,line width=.7pt},
  cutedge/.style={draw=FigureRed,line width=1.5pt},
  lab/.style={font=\figurefont\fontsize{7.5}{9}\selectfont,align=center},
  note/.style={font=\FigureNoteFont\fontsize{7.5}{9}\selectfont,text=FigureMuted,align=left}]
\begin{scope}[shift={(0,0)}]
  \node[vC] (v1) at (4,0) {$1$};
  \node[vC] (v2) at (4,9) {$2$};
  \node[vC] (v3) at (10.5,4.5) {$3$};
  \node[vC] (v4) at (18.5,4.5) {$4$};
  \node[vC] (v5) at (25,9) {$5$};
  \node[vC] (v6) at (25,0) {$6$};
  \node[vN] (v7) at (-2.5,14) {$7$};
  \draw[edge] (v1)--(v2);
  \draw[edge] (v2)--(v3);
  \draw[edge] (v1)--(v3);
  \draw[edge] (v4)--(v5);
  \draw[edge] (v5)--(v6);
  \draw[edge] (v4)--(v6);
  \draw[edge] (v3)--(v4);
  \draw[cutedge] (v2)--(v7);
  \node[lab,anchor=north] at (11.25,-4) {(1) 移入$7$，割权$0.6$\\RatioCut${}=0.7$};
\end{scope}
\begin{scope}[shift={(36,0)}]
  \node[vC] (v1) at (4,0) {$1$};
  \node[vN] (v2) at (4,9) {$2$};
  \node[vC] (v3) at (10.5,4.5) {$3$};
  \node[vC] (v4) at (18.5,4.5) {$4$};
  \node[vC] (v5) at (25,9) {$5$};
  \node[vC] (v6) at (25,0) {$6$};
  \node[vS] (v7) at (-2.5,14) {$7$};
  \draw[cutedge] (v1)--(v2);
  \draw[cutedge] (v2)--(v3);
  \draw[edge] (v1)--(v3);
  \draw[edge] (v4)--(v5);
  \draw[edge] (v5)--(v6);
  \draw[edge] (v4)--(v6);
  \draw[edge] (v3)--(v4);
  \draw[edge] (v2)--(v7);
  \node[lab,anchor=north] at (11.25,-4) {(2) 移入$2$，割权$2$\\RatioCut${}=1.4$};
\end{scope}
\begin{scope}[shift={(72,0)}]
  \node[vN] (v1) at (4,0) {$1$};
  \node[vS] (v2) at (4,9) {$2$};
  \node[vC] (v3) at (10.5,4.5) {$3$};
  \node[vC] (v4) at (18.5,4.5) {$4$};
  \node[vC] (v5) at (25,9) {$5$};
  \node[vC] (v6) at (25,0) {$6$};
  \node[vS] (v7) at (-2.5,14) {$7$};
  \draw[edge] (v1)--(v2);
  \draw[cutedge] (v2)--(v3);
  \draw[cutedge] (v1)--(v3);
  \draw[edge] (v4)--(v5);
  \draw[edge] (v5)--(v6);
  \draw[edge] (v4)--(v6);
  \draw[edge] (v3)--(v4);
  \draw[edge] (v2)--(v7);
  \node[lab,anchor=north] at (11.25,-4) {(3) 移入$1$，割权$2$\\RatioCut${}=\tfrac{7}{6}\approx1.17$};
\end{scope}
\begin{scope}[shift={(0,-33)}]
  \node[vS] (v1) at (4,0) {$1$};
  \node[vS] (v2) at (4,9) {$2$};
  \node[vN] (v3) at (10.5,4.5) {$3$};
  \node[vC] (v4) at (18.5,4.5) {$4$};
  \node[vC] (v5) at (25,9) {$5$};
  \node[vC] (v6) at (25,0) {$6$};
  \node[vS] (v7) at (-2.5,14) {$7$};
  \draw[edge] (v1)--(v2);
  \draw[edge] (v2)--(v3);
  \draw[edge] (v1)--(v3);
  \draw[edge] (v4)--(v5);
  \draw[edge] (v5)--(v6);
  \draw[edge] (v4)--(v6);
  \draw[cutedge] (v3)--(v4);
  \draw[edge] (v2)--(v7);
  \node[lab,anchor=north] at (11.25,-4) {(4) 移入$3$，割权$1$\\\color{FigureRed}RatioCut${}=\tfrac{7}{12}\approx0.58$};
\end{scope}
\begin{scope}[shift={(36,-33)}]
  \node[vS] (v1) at (4,0) {$1$};
  \node[vS] (v2) at (4,9) {$2$};
  \node[vS] (v3) at (10.5,4.5) {$3$};
  \node[vN] (v4) at (18.5,4.5) {$4$};
  \node[vC] (v5) at (25,9) {$5$};
  \node[vC] (v6) at (25,0) {$6$};
  \node[vS] (v7) at (-2.5,14) {$7$};
  \draw[edge] (v1)--(v2);
  \draw[edge] (v2)--(v3);
  \draw[edge] (v1)--(v3);
  \draw[cutedge] (v4)--(v5);
  \draw[edge] (v5)--(v6);
  \draw[cutedge] (v4)--(v6);
  \draw[edge] (v3)--(v4);
  \draw[edge] (v2)--(v7);
  \node[lab,anchor=north] at (11.25,-4) {(5) 移入$4$，割权$2$\\RatioCut${}=1.4$};
\end{scope}
\begin{scope}[shift={(69,-19)}]
  \node[vS,minimum size=3.6mm] at (0,0) {}; \node[note,anchor=west] at (2.8,0) {属于$S$};
  \node[vC,minimum size=3.6mm] at (0,-5) {}; \node[note,anchor=west] at (2.8,-5) {属于$\bar S$};
  \node[vN,minimum size=3.6mm] at (0,-10) {}; \node[note,anchor=west] at (2.8,-10) {本步移入};
  \draw[cutedge] (-1.8,-15)--(1.8,-15); \node[note,anchor=west] at (2.8,-15) {跨组边};
\end{scope}
\end{tikzpicture}
